\appendix

\chapter{附录}
\label{app:basic-notation-topology}

众所周知, 空间 \(\mathbb R^d\) 是实数域 \(\mathbb R\) 上的一个向量空间. 本书前面各章中的积分理论, 都是在这个空间及其开子集上建立的. \(\mathbb R^d\) 上的每一个线性映射都可写成
\[
x\longmapsto Ax,
\]
其中 \(A\) 是一个 \(d\times d\) 实矩阵, 而 \(x\) 可以看成一个 \(d\times 1\) 列矩阵. 所有可逆 \(d\times d\) 实矩阵组成的群记为
\[
\operatorname{GL}(d,\mathbb R).
\]

若 \(x=(x_1,\ldots,x_d)\in\mathbb R^d\), 我们总用
\[
\|x\|=(x_1^2+\cdots+x_d^2)^{1/2}
\]
表示 \(x\) 的 Euclidean 范数.

\section{拓扑记号}
\label{app:topological-notation}

回忆 \(\mathbb R^d\) 中的开集定义. 若 \(O\subset\mathbb R^d\), 且对每个 \(x\in O\), 都存在 \(\delta>0\), 使得
\[
\{y:\|y-x\|<\delta\}\subset O,
\]
则称 \(O\) 为开集. 开集的任意并集以及有限交集仍是开集.

因此, 包含在给定集合 \(M\subset\mathbb R^d\) 中的所有开集的并仍是开集. 这个开集称为 \(M\) 的内部, 记作
\[
M^\circ.
\]
于是 \(x\in M^\circ\) 的充分必要条件是: 存在以 \(x\) 为中心的开球包含在 \(M\) 中.

若集合 \(F\subset\mathbb R^d\) 的余集
\[
F^c=\mathbb R^d\setminus F
\]
是开集, 则称 \(F\) 为闭集. 闭集的任意交集以及有限并集仍是闭集. 包含给定集合 \(M\) 的所有闭集的交称为 \(M\) 的闭包, 记作
\[
\overline M.
\]
也就是说, \(\overline M\) 是包含 \(M\) 的最小闭集. 点 \(x\in\overline M\) 的充分必要条件是: 每个以 \(x\) 为中心的开球都与 \(M\) 相交.

集合
\[
\partial M=\overline M\cap \overline{M^c}
\]
称为 \(M\) 的边界. 因此, \(x\in\partial M\) 当且仅当任意以 \(x\) 为中心的开球既与 \(M\) 相交, 又与 \(M^c\) 相交.

\begin{defn}[紧集]
	\label{def:compact-set-appendix}
	\(\mathbb R^d\) 中的一个闭且有界的集合称为紧集.
\end{defn}

根据 Heine--Borel 定理, 紧集还可以用开覆盖来刻画: 若 \(K\) 为紧集, 则覆盖 \(K\) 的任意一族开集必有有限子族仍覆盖 \(K\).

\begin{exer}
	\label{exer:closed-sequential-characterization}
	证明:
	\begin{enumerate}
		\item 集合 \(F\subset\mathbb R^d\) 为闭集的充分必要条件是: \(F\) 中任意收敛序列的极限仍属于 \(F\).
		\item 集合 \(K\subset\mathbb R^d\) 为紧集的充分必要条件是: \(K\) 中任意序列都有一个收敛子序列, 且该子序列的极限仍属于 \(K\).
	\end{enumerate}
\end{exer}

\section{函数和集合的基本记号}
\label{app:function-set-notation}

对于函数
\[
f:\mathbb R^d\longrightarrow\mathbb R,
\]
其支集 \(\operatorname{supp}f\) 定义为集合
\[
\{x:f(x)\neq0\}
\]
的闭包. 等价地, \((\operatorname{supp}f)^c\) 是使 \(f(x)=0\) 的点集的内部. 因为支集总是闭集, 所以 \(\operatorname{supp}f\) 为紧集的充分必要条件是: \(f\) 在某个有界集之外恒等于零.

我们把从 \(\mathbb R^d\) 到 \(\mathbb R\) 的所有连续函数组成的集合记为
\[
C=C(\mathbb R^d)=C(\mathbb R^d,\mathbb R),
\]
把其中具有紧致支集的函数组成的集合记为
\[
C_0=C_0(\mathbb R^d).
\]

对于任意函数列 \(f_n\), 可逐点定义
\[
|f|,\qquad f_1+f_2,\qquad f_1f_2,\qquad
\max(f_1,f_2),\qquad \min(f_1,f_2),
\]
以及
\[
\lim f_n,\qquad \sup_n f_n,\qquad \inf_n f_n.
\]
例如
\[
\bigl(\max(f_1,f_2)\bigr)(x)=\max\{f_1(x),f_2(x)\},
\]
而
\[
\bigl(\sup_n f_n\bigr)(x)=\sup_n f_n(x),
\qquad x\in\mathbb R^d.
\]

我们记
\[
f^+=\max(f,0)\geq0,
\qquad
f^-=\max(-f,0)\geq0.
\]
于是
\[
f=f^+-f^-,
\qquad
|f|=f^++f^-.
\]

若 \(E\subset\mathbb R^d\), 其特征函数记为
\[
\chi_E(x)=
\begin{cases}
	1, & x\in E,\\
	0, & x\in E^c=\mathbb R^d\setminus E.
\end{cases}
\]
于是
\[
\chi_{E^c}=1-\chi_E,
\]
并且对任意集合 \(E_1,E_2\subset\mathbb R^d\), 有
\[
\chi_{E_1\cup E_2}
=
\max(\chi_{E_1},\chi_{E_2})
=
\chi_{E_1}+\chi_{E_2}-\chi_{E_1}\chi_{E_2},
\]
\[
\chi_{E_1\cap E_2}
=
\min(\chi_{E_1},\chi_{E_2})
=
\chi_{E_1}\chi_{E_2},
\]
以及
\[
\chi_{E_1\triangle E_2}
=
|\chi_{E_1}-\chi_{E_2}|
=
\chi_{E_1}+\chi_{E_2}-2\chi_{E_1}\chi_{E_2}.
\]
这里
\[
E_1\triangle E_2=(E_1\cup E_2)\setminus(E_1\cap E_2)
\]
称为 \(E_1\) 和 \(E_2\) 的对称差, 即由属于 \(E_1\) 与 \(E_2\) 中恰好一个集合的点组成.

\begin{exer}
	\label{exer:indicator-continuity-boundary}
	证明: \(\chi_E\) 在点 \(x\) 连续的充分必要条件是
	\[
	x\notin\partial E.
	\]
\end{exer}

\section{极限、上下极限和上确界}
\label{app:limsup-sup-notation}

对于序列 \((a_n)\), 其中每个 \(a_n\) 可以是实数, 也可以是 \(+\infty\) 或 \(-\infty\), 若
\[
a_n\leq a_{n+1}
\qquad
(\text{或 }a_n\geq a_{n+1})
\]
且
\[
\lim_{n\to\infty}a_n=a,
\]
则分别用
\[
a_n\uparrow a,
\qquad
a_n\downarrow a
\]
表示.

对于任意序列 \((a_n)\), 优于它的最小下降序列是
\[
\left(\sup_{k\geq n}a_k\right)_{n=1}^{\infty}.
\]
因此定义
\[
\overline{\lim}_{n\to\infty}a_n
=
\lim_{n\to\infty}\sup_{k\geq n}a_k
=
\inf_n\sup_{k\geq n}a_k.
\]
类似地, 定义
\[
\underline{\lim}_{n\to\infty}a_n
=
\lim_{n\to\infty}\inf_{k\geq n}a_k
=
\sup_n\inf_{k\geq n}a_k.
\]
显然
\[
\underline{\lim}_{n\to\infty}a_n
\leq
\overline{\lim}_{n\to\infty}a_n,
\]
且等号成立当且仅当 \(\lim_{n\to\infty}a_n\) 存在.

设
\[
f:M\longrightarrow\mathbb R
\]
是任意函数. 我们分别用
\[
\sup_M f,
\qquad
\inf_M f
\]
表示值集 \(f(M)\) 的上确界和下确界. 若 \(f\) 上方无界, 则约定
\[
\sup_M f=+\infty.
\]
显然
\[
\inf_M f=-\sup_M(-f),
\]
并且
\[
\sup_M(f+g)
\leq
\sup_M f+
\sup_M g.
\]
取 \(M=\{k:k\geq n\}\), 便可得到: 若 \((a_n)\) 和 \((b_n)\) 为上方有界序列, 则
\[
\overline{\lim}_{n\to\infty}(a_n+b_n)
\leq
\overline{\lim}_{n\to\infty}a_n
+
\overline{\lim}_{n\to\infty}b_n.
\]
类似地, 若二者下方有界, 则
\[
\underline{\lim}_{n\to\infty}(a_n+b_n)
\geq
\underline{\lim}_{n\to\infty}a_n
+
\underline{\lim}_{n\to\infty}b_n.
\]

若 \(f\) 和 \(g\) 上方有界, 则由不等式
\[
f\leq g+|f-g|,
\qquad
g\leq f+|f-g|
\]
可得
\[
\left|\sup_M f-\sup_M g\right|
\leq
\sup_M|f-g|.
\]
当 \(M=\{1,2\}\) 时, 这给出
\[
\left|\max(a_1,a_2)-\max(b_1,b_2)\right|
\leq
\max\{|a_1-b_1|,|a_2-b_2|\}.
\]

\section{距离函数和截断函数}
\label{app:distance-function-cutoff}

对于 \(\mathbb R^d\) 的非空子集 \(M\), 定义点 \(x\) 到 \(M\) 的距离为
\[
d(x,M)=\inf_{y\in M}\|x-y\|
=
-\sup_{y\in M}\bigl(-\|x-y\|\bigr).
\]
注意
\[
d(x,M)=0
\]
的充分必要条件是
\[
x\in\overline M.
\]
其次, 由
\begin{align*}
	\left|d(x_1,M)-d(x_2,M)\right|
	&\leq
	\sup_{y\in M}\left|-\|x_1-y\|+\|x_2-y\|\right|\\
	&\leq \|x_1-x_2\|
\end{align*}
可知, 函数
\[
x\longmapsto d(x,M)
\]
是连续的, 事实上是一致连续的.

\begin{exer}
	\label{exer:distance-cutoff-open-closed}
	设 \(F\) 是开集 \(O\) 的一个闭子集. 证明函数
	\[
	f(x)=
	\frac{d(x,O^c)}{d(x,F)+d(x,O^c)},
	\qquad x\in\mathbb R^d,
	\]
	是连续的, 并且
	\[
	0\leq f\leq1,
	\qquad
	f(x)=1\quad(x\in F),
	\qquad
	\operatorname{supp}f\subset \overline O.
	\]
	其次证明: 若 \(K\subset O\) 是紧集, 则存在 \(f\in C_0(\mathbb R^d)\), 使得
	\[
	0\leq f\leq1,
	\qquad
	f=1\text{ 于 }K\text{ 的某个邻域内},
	\qquad
	\operatorname{supp}f\subset O.
	\]
\end{exer}

\chapter{综合练习}
\label{chap:comprehensive-exercises}

\begin{exer}
	设 \(f\) 是 \(\R^d\) 上的一个函数, 满足
	\[
	f(x+l)=f(x),\qquad x,l\in\R^d,
	\]
	其中 \(l\) 具有整数坐标. 证明: 若 \(f\) 在
	\[
	I=\{x:0\leq x_j\leq 1,\ j=1,\ldots,d\}
	\]
	上 Riemann 可积, 则
	\[
	\lim_{T\to\infty}\frac1T\int_0^T
	f(t\alpha_1,\ldots,t\alpha_d)\,\dd t
	=
	\int_I f(x)\,\dd x,
	\]
	其中 \(\alpha_1,\ldots,\alpha_d\) 在有理数域上线性无关. 也就是说, 若
	\[
	r_1\alpha_1+\cdots+r_d\alpha_d=0,
	\qquad r_j\in\mathbb Q,
	\]
	则必有 \(r_1=\cdots=r_d=0\).
\end{exer}

\begin{exer}
	设 \(f\) 是开集 \(O\subset\R^d\) 上的二次连续可微函数, 并令
	\[
	f_h(x)=f(x)-\langle x,h\rangle .
	\]
	证明: 对除一个零集外的所有 \(h\in\R^d\), 矩阵
	\[
	\left(\frac{\partial^2 f}{\partial x_j\partial x_k}\right)
	\]
	在 \(O\) 中使
	\[
	\frac{\partial f_h}{\partial x_j}=0,
	\qquad j=1,\ldots,d,
	\]
	的任意一点上都是可逆的.
\end{exer}

\begin{exer}
	在实数轴上定义函数 \(f\): 当 \(x\) 为无理数时 \(f(x)=0\), 并且
	\[
	f(0)=0,
	\qquad
	f(x)=\frac1q\quad\text{当 }x=\frac pq,
	\]
	其中 \(p,q\) 互素且 \(q>0\). 证明 \(f\) 是上方半连续函数.
\end{exer}

\begin{exer}
	从闭区间 \([0,1]\) 的中央删去一个长度为 \(1/p_1\) 的开区间. 再从余下各个闭子区间的中央同时删去一个开区间, 其长度是该子区间长度的 \(1/p_2\). 如此重复进行, 作到第 \(n\) 步得到闭集 \(F_n\). 证明
	\[
	m\left(\bigcap_{n=1}^{\infty}F_n\right)
	=
	\prod_{n=1}^{\infty}\left(1-\frac1{p_n}\right).
	\]
\end{exer}

\begin{exer}
	求所有在十进制记数法的表示中不出现数字 \(7\) 的实数所成集合的测度.
\end{exer}

\begin{exer}
	设 \(E\) 是一个实数集合. 当 \(E\) 中的数用十进制记数法表示时, 任意一个有限的数字序列都在 \(E\) 中出现. 证明 \(E^c\) 是一个零集.
\end{exer}

\begin{exer}
	设 \(E_1,E_2,\ldots\) 是 \(\R^d\) 的子集, 并且
	\[
	\sum_{j=1}^{\infty}m^*(E_j)<+\infty.
	\]
	证明集合
	\[
	\{x:x\in E_j\text{ 对无限多个 }j\text{ 成立}\}
	\]
	是零集.
\end{exer}

\begin{exer}
	设 \((a_n)_{n=1}^{\infty}\) 为实数序列, 并假定存在常数 \(K\), 使得任一长度为 \(1\) 的区间都至多包含 \(K\) 个 \(a_n\). 证明: 若 \(f\in L^1(\R)\), 则级数
	\[
	\sum_{n=1}^{\infty} f(x+a_n)
	\]
	对几乎所有的 \(x\) 绝对收敛.
\end{exer}

\begin{exer}
	证明: 若 \(f\in L^1(\R)\), 则级数
	\[
	\sum_{n=1}^{\infty}f(nx)
	\]
	对几乎所有的 \(x\) 绝对收敛.
\end{exer}

\begin{exer}
	设 \((a_n)_{n=1}^{\infty}\) 和 \((b_n)_{n=1}^{\infty}\) 为实数序列, 且 \(b_n\geq0\). 证明级数
	\[
	\sum_{n=1}^{\infty} b_n |x-a_n|^{-\alpha}
	\]
	当
	\[
	\sum_{n=1}^{\infty} b_n^{1/\alpha}<+\infty,
	\qquad \alpha>1,
	\]
	时对几乎所有的 \(x\) 收敛.
\end{exer}

\begin{exer}
	对于一个实数 \(a\), 令 \(E_a\) 为 \(\R^d\) 中所有满足下述条件的 \(x\) 之集合: 不等式
	\[
	|\langle l,x\rangle|
	=|l_1x_1+\cdots+l_dx_d|
	\geq \|l\|^{-a}
	\]
	对除有限个之外的所有整数向量 \(l=(l_1,\ldots,l_d)\) 成立. 证明: 若 \(a<d-1\), 则 \(E_a^c\) 是一个零集. 并且证明: 当 \(d>1\) 时, 对每个 \(x\in\R^d\), 有
	\[
	\liminf_{\|l\|\to\infty}
	\|l\|^{d-1}|\langle l,x\rangle|
	\leq C_d\|x\|,
	\]
	其中 \(C_d\) 是只依赖于 \(d\) 的常数.
\end{exer}

\begin{exer}
	设 \(E\) 为实数轴的一个子集. 用 \(M\) 表示所有与 \(E\) 中无穷多个点模 \(1\) 相等的点所组成的集合. 证明: 若 \(E\) 可测, 则 \(M\) 可测; 若 \(E\) 有有限外测度, 则 \(M\) 是一个零集.
\end{exer}

\begin{exer}
	设 \(E\) 为实数轴上的一个零集. 证明: 存在实数 \(a\), 使得对任意有理数 \(r\), 都有
	\[
	a+r\notin E.
	\]
\end{exer}

\begin{exer}
	设 \((E_k)_{k=1}^{\infty}\) 为 \(\R^d\) 中可测集合的一个序列. 对 \(i=0,1,2,\ldots\), 用 \(A_i\) 表示那些恰好属于 \(i\) 个集合 \(E_k\) 的点所组成的集合. 再令 \(A_\infty\) 表示属于无穷多个 \(E_k\) 的点所组成的集合. 证明 \(A_i\) 均可测, 并且在约定 \(+\infty\cdot m(A_\infty)\) 的意义下有
	\[
	\sum_{k=1}^{\infty}m(E_k)
	=
	\sum_{i=1}^{\infty} i m(A_i)+(+\infty)m(A_\infty).
	\]
	特别地, 若 \(m(A_\infty)=0\), 则
	\[
	\sum_{k=1}^{\infty}m(E_k)
	=
	\sum_{i=1}^{\infty} i m(A_i).
	\]
\end{exer}

\begin{exer}
	设 \(E\) 是 \((0,1)\) 中的一个可测集, 且 \(m(E)>1/2\). 证明存在 \(x,y\in E\), 使得
	\[
	x+y=1.
	\]
\end{exer}

\begin{exer}
	证明: 若实数轴上 \(f\in\mathcal I^+\), 则
	\[
	\int_{\R}^{*} f(x)\,\dd x
	\leq
	\liminf_{n\to\infty}\frac1n\sum_{k=-n^2}^{n^2}f\left(\frac{k}{n}\right).
	\]
\end{exer}

\begin{exer}
	设 \(f\in L^p(\R^d)\), \(g\in L^q(\R^d)\), 其中
	\[
	1\leq p,q\leq +\infty,
	\qquad
	\frac1p+\frac1q=1.
	\]
	证明函数
	\[
	x\longmapsto \int_{\R^d} f(x-y)g(y)\,\dd y
	\]
	是 \(x\) 的连续函数.
\end{exer}

\begin{exer}
	设 \(A,B\subset\R^d\) 为可测集合, 并假定对每个 \(x\in\R^d\), 存在一个零集 \(N_x\), 使得
	\[
	A+x\subset B\cup N_x.
	\]
	证明: 若 \(A\) 不是零集, 则 \(B^c\) 是零集.
\end{exer}

\begin{exer}
	假定 \(A\) 是一个可测且具有正测度的实数集合. 证明
	\[
	\{x-y:x,y\in A\}
	\]
	包含原点的一个邻域.
\end{exer}

\begin{exer}
	证明: 若实数轴上 \(f\in L^1\), 则
	\[
	\lim_{n\to\infty}
	\sum_{k=-n^2}^{n^2}
	\left|
	\int_{k/n}^{(k+1)/n}f(x)\,\dd x
	\right|
	=
	\int_{-\infty}^{+\infty}|f(x)|\,\dd x.
	\]
\end{exer}

\begin{exer}
	设 \(f\in L^1(\R)\), \(g\in L^\infty(\R)\), 并且 \(g\) 是周期为 \(T\) 的函数. 证明
	\[
	\lim_{n\to\infty}\int_{\R} f(x)g(nx)\,\dd x
	=
	\left(\int_{\R}f(x)\,\dd x\right)
	\left(\frac1T\int_0^T g(x)\,\dd x\right).
	\]
\end{exer}

\begin{exer}
	证明: 若级数
	\[
	\sum_{n=1}^{\infty}a_n\sin nx
	\]
	在一个正测度集合中的所有 \(x\) 处绝对收敛, 则
	\[
	\sum_{n=1}^{\infty}|a_n|<+\infty.
	\]
\end{exer}

\begin{exer}
	证明: 若 \(f\in L^1(\R)\), 则
	\[
	\lim_{T\to\infty}
	\int_{-\infty}^{+\infty}
	\left|
	\frac1{2T}\int_{-T}^{T} f(s+t)\,\dd t
	\right|\dd s
	=
	\left|\int_{-\infty}^{+\infty}f(t)\,\dd t\right|.
	\]
\end{exer}

\begin{exer}
	设 \(f_1,f_2,\ldots\) 是区间 \((0,1)\) 上的可测函数, 并且
	\[
	|f_n(x)|\leq1.
	\]
	证明: 若
	\[
	\int_0^y f_n(x)\,\dd x\to0,
	\qquad n\to\infty,
	\]
	对所有 \(y\in(0,1)\) 成立, 则对所有 \(g\in L^1(0,1)\), 有
	\[
	\int_0^1 g(x)f_n(x)\,\dd x\to0,
	\qquad n\to\infty.
	\]
\end{exer}

\begin{exer}
	设 \(f\in L^1\). 证明: 对可测集 \(E\), 当 \(m(E)\to0\) 时,
	\[
	\int_E f(x)\,\dd x\to0.
	\]
	更强地, 有
	\[
	\int_E |f(x)|\,\dd x\to0.
	\]
\end{exer}

\begin{exer}
	设 \(0\leq f_n\in L^1\). 假定对几乎所有的 \(x\), 有
	\[
	f_n(x)\to f(x),
	\]
	并且
	\[
	\int f_n(x)\,\dd x\to A,
	\qquad n\to\infty.
	\]
	证明 \(f\in L^1\), 并且
	\[
	\int |f_n-f|\,\dd x\to A-\int f\,\dd x,
	\qquad n\to\infty.
	\]
\end{exer}

\begin{exer}
	设 \(f,f_1,f_2,\ldots\in L^p\), 其中 \(1\leq p<+\infty\). 假定对几乎所有的 \(x\), 有
	\[
	f_n(x)\to f(x),
	\]
	并且
	\[
	\|f_n\|_p\to \|f\|_p.
	\]
	证明
	\[
	\|f_n-f\|_p\to0.
	\]
\end{exer}

\begin{exer}
	设 \(f\in L^p\), \(g_n\in L^q\), \(n=1,2,\ldots\), 其中
	\[
	1<p<+\infty,
	\qquad
	\frac1p+\frac1q=1.
	\]
	证明: 若 \(\|g_n\|_q\leq C\), 且 \(g_n(x)\to g(x)\) 几乎处处成立, 则 \(g\in L^q\), 并且
	\[
	\int f g_n\,\dd x\to \int f g\,\dd x.
	\]
\end{exer}

\begin{exer}
	设 \(f_n\in L^p\), \(g_n\in L^q\), 其中
	\[
	1<p<+\infty,
	\qquad
	\frac1p+\frac1q=1.
	\]
	证明: 若
	\[
	|f_n|\leq F\in L^p\quad\text{对所有 }n,
	\qquad
	\|g_n\|_q\leq C,
	\]
	且
	\[
	f_n(x)\to f(x),
	\qquad
	g_n(x)\to g(x)
	\]
	几乎处处成立, 则 \(f\in L^p\), \(g\in L^q\),
	\[
	\|f-f_n\|_p\to0,
	\]
	并且
	\[
	\int f_n g_n\,\dd x\to \int f g\,\dd x.
	\]
\end{exer}

\begin{exer}
	设 \(f_n\) 为非负可测函数序列, 且
	\[
	\int f_n(x)\,\dd x=1
	\]
	对每个 \(n\) 成立, 并且 \(f_n(x)\to0\) 对所有 \(x\) 成立. 试问函数
	\[
	f=\sup_n f_n
	\]
	是否一定可测? 是否一定可积?
\end{exer}

\begin{exer}
	设 \(f(x,y)\) 对 \((x,y)\in\R^2\) 有定义. 假定 \(f(x,y)\) 对任一固定 \(y\) 是 \(x\) 的连续函数, 并且对任一固定 \(x\) 是 \(y\) 的可测函数. 证明 \(f\) 是 \(\R^2\) 上的可测函数.
\end{exer}

\begin{exer}
	试问 \((0,1)\) 中一个开子集的边界是否必定是零测集?
\end{exer}

\begin{exer}
	设 \(E\) 是一个实数集合, 并且存在常数 \(k<1\), 使得对任意区间 \(I\), 有
	\[
	m^*(E\cap I)\leq k m(I).
	\]
	证明 \(E\) 是零集.
\end{exer}

\begin{exer}
	构造一个可测集合 \(E\subset\R^d\), 使得对任一测度为有限值的非空开集 \(O\), 都有
	\[
	0<m(E\cap O)<m(O).
	\]
\end{exer}

\begin{exer}
	设 \(E\) 是区间 \(I\subset\R^d\) 的一个子集, 并令 \(0<\varepsilon<1\). 若 \(I_1,\ldots,I_N\) 是区间 \(I\) 的一个剖分, 用 \(I'\) 表示所有满足
	\[
	m^*(I_j\cap E)>\varepsilon m(I_j)
	\]
	的 \(I_j\) 的并. 证明当剖分的细度趋于零时,
	\[
	m(I')\to m^*(E).
	\]
	再用同样方法定义 \(I''\), 但将 \(E\) 替换为 \(I\setminus E\). 证明 \(E\) 可测的充分必要条件是
	\[
	m(I'\cap I'')\to0
	\]
	随剖分细度趋于零成立.
\end{exer}

\begin{exer}
	设 \(f_n\) 是在 \((-\infty,+\infty)\) 中取值的可测函数序列. 证明: 若对每个 \(\varepsilon>0\), 有
	\[
	m\{x:|f_m(x)-f_n(x)|\geq\varepsilon\}\to0,
	\qquad n,m\to\infty,
	\]
	则可以选出一个子序列, 使它对几乎所有的 \(x\) 收敛.
\end{exer}

\begin{exer}
	设非负可测函数 \(f\) 和 \(g\) 定义在可测集 \(E\subset\R^d\) 上. 令
	\[
	E_y=\{x\in E:g(x)\geq y\},
	\qquad
	h(y)=\int_{E_y}f(x)\,\dd x.
	\]
	证明
	\[
	\int_E f(x)g(x)\,\dd x
	=
	\int_0^{\infty}h(y)\,\dd y.
	\]
\end{exer}

\begin{exer}
	假定 \(f\) 在区间 \([0,1]\) 上连续, 而 \(g\) 是 \([0,1]\) 上的可测函数, 且
	\[
	0\leq g(x)\leq1.
	\]
	证明极限
	\[
	\lim_{n\to\infty}\int_0^1 f\bigl(g(x)^n\bigr)\,\dd x
	\]
	存在, 并给出其值.
\end{exer}

\begin{exer}
	设 \(f\) 在实数轴上可积. 证明对几乎所有的 \(a\), 有
	\[
	\lim_{n\to\infty}\int_0^1 f(x^n+a)\,\dd x=f(a).
	\]
\end{exer}

\begin{exer}
	设 \(E\) 是 \(\R\) 的子集. 证明
	\[
	m^*(E)
	=
	\inf_{f\in K}
	\left\{
	\lim_{x\to+\infty}\bigl(f(x)-f(-x)\bigr)
	\right\},
	\]
	其中 \(K\) 是所有在 \(\R\) 上递增, 并且在 \(E\) 的每一点处微商存在且不小于 \(1\) 的函数 \(f\) 所组成的集合.
\end{exer}

\begin{exer}
	设 \(f\) 是 \((0,1)\) 上的一个增函数. 证明
	\[
	\int_0^1 f'(x)\,\dd x
	\]
	最大不超过 \(f\) 的值域的测度.
\end{exer}

\chapter{综合练习的提示}
\label{chap:comprehensive-exercise-hints}

\begin{enumerate}
	\item 根据 Weierstrass 逼近定理, 一个连续周期函数可以用三角多项式
	\[
	\sum c_l \exp(2\pi i\langle l,x\rangle)
	\]
	一致逼近, 其中 \(l\) 的坐标为整数. 修改式 \((1.1.5)\), 使其中的
	\(h_1\) 和 \(h_2\) 为三角多项式. 再思考
	\[
	f(x)=\exp(2\pi i\langle l,x\rangle)
	\]
	这种特征函数情形下结论的含义.
	
	\item 对映射 \(f'\) 应用 Morse--Sard 定理. 注意, 若使用 Lebesgue 零集, 则不必限制在紧集 \(K\) 上.
	
	\item 对任意 \(a>0\), 集合
	\[
	\{x:f(x)\geq a\}
	\]
	是离散集, 因而是闭集. 其次, 当 \(a\leq0\) 时, 该集合为全体实数集.
	
	\item 使用 Beppo Levi 定理.
	
	\item 这只是练习 4 的一个变体. 这里每次删去区间长度的 \(1/10\), 只是删去的不一定是正中间的那个子区间. 所得集合的测度为零.
	
	\item 当只考虑一个固定的有限数字串时, 问题与练习 5 基本相同. 最后注意, 有限数字串只有可数多个.
	
	\item 考虑特征函数之和.
	
	\item 与定理 \ref{thm:riesz-fischer-l1} 前面的练习比较. 证明
	\[
	\sum_{n=1}^{\infty}|f(x+a_n)|
	\]
	是局部可积的.
	
	\item 考虑
	\[
	\int g(x)\sum_{n=1}^{\infty}|f(nx)|\,\dd x,
	\]
	其中 \(g\in C_0^+\), 且 \(g\) 在靠近 \(0\) 的地方恒等于零.
	
	\item 如果函数 \(|x-a_n|^{-\alpha}\) 是可积的, 那么这基本上是练习 8. 若单独考虑满足
	\[
	|x-a_n|\leq b_n^{1/\alpha}
	\]
	对无穷多个 \(n\) 成立的那些点 \(x\), 则先证明这些点构成一个零集, 然后排除这些奇性并像练习 8 那样进行证明.
	
	\item 估计所有满足
	\[
	\|x\|\leq M,
	\qquad
	|l_1x_1+\cdots+l_dx_d|\leq \|l\|^{-a}
	\]
	的 \(x\) 的测度. 判断
	\[
	\sum_{l\ne0}\|l\|^{-b}
	\]
	在什么条件下收敛.
	
	若不用积分理论直接证明第二个论断会比较复杂. 当然, 若存在 \(l\ne0\) 使得
	\[
	\langle l,x\rangle=0,
	\]
	则结论平凡. 在其他情形, 考虑所有纯量积
	\[
	\langle l,x\rangle,
	\qquad \|l\|\leq R.
	\]
	它们是互不相同的, 并且都位于区间
	\[
	(-R\|x\|,R\|x\|)
	\]
	之内. 因为这样的整数向量 \(l\) 的个数至少为
	\[
	C(R-\sqrt d)^d,
	\]
	其中 \(C\) 是单位球的体积, 所以可以找到 \(l_1(R)\) 和 \(l_2(R)\), 使得
	\[
	0<
	|\langle x,l_1(R)\rangle-\langle x,l_2(R)\rangle|
	\leq
	2R\|x\|
	\bigl(C(R-\sqrt d)^d-1\bigr)^{-1}\to0.
	\]
	于是
	\[
	\liminf_{R\to\infty}
	\|l(R)\|^{d-1}|\langle x,l(R)\rangle|
	\leq
	2^dC^{-1}\|x\|,
	\]
	其中
	\[
	l(R)=l_1(R)-l_2(R),
	\qquad \|l(R)\|\to\infty.
	\]
	
	\item 若 \(f\) 是 \(E\) 或 \(E^c\) 中某个可积集合的特征函数, 则有练习 8 的论断.
	
	\item 可数多个零集之并仍是零集.
	
	\item 引进特征函数.
	
	\item 否则 \(E\) 和
	\[
	E_1=\{x:1-x\in E\}
	\]
	就应当是分离的. 这样它们的并集包含于 \((0,1)\), 但测度大于 \(1\), 矛盾.
	
	\item 若 \(f\) 为连续函数, 则等号成立. 对一般的 \(f\in\mathcal I^+\), 用连续函数从下方近似 \(f\).
	
	\item 由 H\"older 不等式有
	\[
	\left|
	\int f(x-y)g(y)\,\dd y
	\right|
	\leq
	\|f\|_p\|g\|_q.
	\]
	再按 \(L^p\) 或 \(L^q\) 范数, 用连续函数近似 \(f\) 或 \(g\).
	
	\item 引进 \(A\) 和 \(B^c\) 的可积子集的特征函数, 并利用 Fubini 定理后面的练习.
	
	\item 若 \(f\) 为 \(E\) 的特征函数, 则所考虑的集合包含所有满足
	\[
	\int f(x+y)f(y)\,\dd y\neq0
	\]
	的点 \(x\). 为什么? 可以假定 \(E\) 具有有限测度. 由练习 17, 上述积分是 \(x\) 的连续函数, 并且在原点处不等于零.
	
	\item 当 \(f\in C_0\) 时结论显然. 一般情形根据 \(L^1\) 的定义作近似.
	
	\item 可以先假定 \(g\geq0\). 若 \(f\in C_0\), 结论显然成立: 将积分区间分解成 \(nx\) 落在 \(g\) 的若干周期中的部分, 然后应用积分中值定理. 对一般 \(f\in L^1\), 用 \(C_0\) 函数逼近.
	
	当 \(g\) 为正弦函数或余弦函数时, 这个特殊情形称为 Riemann--Lebesgue 引理.
	
	\item 各项绝对值之和在某个正测度集合上是有界的. 在这个集合上积分, 并利用练习 21.
	
	\item 当 \(f\in C_0\) 时结论容易. 一般情形用 \(C_0\) 函数逼近.
	
	\item 首先考虑阶梯函数 \(g\) 的情形; 这些阶梯函数在 \(L^1\) 中稠密.
	
	\item 用一个 \(C_0\) 中的函数近似 \(f\).
	
	\item 利用 Fatou 引理, 并考虑
	\[
	g_n=\min(f,f_n).
	\]
	注意 \(A=\int f\,\dd x\) 的特殊情形.
	
	\item 根据 Lusin 定理和练习 25, 可以找到紧集 \(K\), 使得
	\[
	\int_{K^c}|f|^p\,\dd x
	\]
	任意小, 同时还使 \(f_n\) 在 \(K\) 上一致收敛到 \(f\). 再利用练习 26.
	
	\item 利用 Lusin 定理、练习 25 和 H\"older 不等式.
	
	\item 利用练习 28.
	
	\item \begin{enumerate}
		\item \(f=\sup_n f_n\) 显然是可测函数.
		\item \(f\) 不一定可积. 否则由 Lebesgue 控制收敛定理可推出
		\[
		\int f_n\,\dd x\to0,
		\]
		这与 \(\int f_n\,\dd x=1\) 矛盾.
	\end{enumerate}
	
	\item 令
	\[
	X_k(x)=\frac jk,
	\qquad
	\frac jk\leq x<\frac{j+1}{k},
	\]
	其中 \(j\) 为整数. 则
	\[
	f(x,y)=\lim_{k\to\infty}f(X_k(x),y),
	\]
	因而 \(f\) 是可测函数.
	
	\item 不一定. 可以由练习 4 得到一个反例.
	
	\item 先把不等式推广到开集 \(I\). 可以假定
	\[
	m^*(E)<+\infty.
	\]
	于是可得
	\[
	m^*(E)\leq k\,m^*(E),
	\]
	从而 \(m^*(E)=0\).
	
	\item 注意, 只需关于一个开集基
	\[
	O_1,O_2,\ldots
	\]
	来作构造即可. 逐次构造 \(E_k\), 使它们关于
	\[
	O_1,\ldots,O_k
	\]
	满足要求, 并且
	\[
	m(E_k\triangle E_{k+1})
	\]
	快速趋于零.
	
	\item 可以假定 \(E\) 是可测集. 先选取有限个区间的并 \(A\), 使得
	\[
	m(E\triangle A)
	\]
	很小. 取 \(\delta<\min(\varepsilon,1-\varepsilon)\). 假定
	\[
	m(I_j\cap(A\triangle E))>\delta m(I_j)
	\]
	不成立. 于是有
	\[
	m(I_j\cap A)>(\varepsilon+\delta)m(I_j)
	\Longrightarrow
	m(I_j\cap E)>\varepsilon m(I_j),
	\]
	以及
	\[
	m(I_j\cap E)>\varepsilon m(I_j)
	\Longrightarrow
	m(I_j\cap A)>(\varepsilon-\delta)m(I_j).
	\]
	这些区间的测度有多大?
	
	\item 参见定理 \ref{thm:riesz-fischer-l1} 的证明.
	
	\item 使用 Fubini 定理. 不要忘记检验可测性.
	
	\item Lebesgue 控制收敛定理给出极限值
	\[
	f(0)m\{x:g(x)<1\}+f(1)m\{x:g(x)=1\}.
	\]
	
	\item 若
	\[
	F(t)=\int_0^t f(a+x)\,\dd x,
	\]
	则
	\[
	\frac{F(t)}{t}\to f(a)
	\]
	对几乎所有的 \(a\) 成立. 可将积分写成
	\[
	\int_0^1 F'(t)\,\dd t^{1/n}
	=
	\frac{F(1)}{n}
	+
	\frac1n\left(1-\frac1n\right)
	\int_0^1 F(t)t^{1/n-2}\,\dd t.
	\]
	先考虑原点附近的积分, 在那里
	\[
	\left|\frac{F(t)}{t}-f(a)\right|<\varepsilon.
	\]
	然后再考虑这个邻域外的积分.
	
	\item 方法一: 令 \(f\) 是包含 \(E\) 的一个开集之特征函数的积分.
	
	方法二: 利用
	\[
	\lim_{x\to+\infty}\bigl(f(x)-f(-x)\bigr)
	\geq
	\lim_{x\to+\infty}\int_{-x}^{x}f'(t)\,\dd t,
	\]
	以及在一个包含 \(E\) 的可测集合上 \(f'\) 存在并且 \(f'\geq1\).
	
	\item 值域的测度等于
	\[
	f(1)-f(0)
	\]
	减去所有不连续点处跳跃度之和. 也就是说, 它等于
	\[
	\int_0^1 f'(x)\,\dd x+\int_0^1 \dd s(x),
	\]
	其中 \(\dd s\) 是 \(\dd f\) 的奇异弥散部分.
\end{enumerate}